\documentclass[runningheads]{llncs}

\usepackage{graphicx}
\usepackage{url}

\begin{document}

\title{MIRAGE-CARMA: Context-Adaptive Retrieval with Metadata Awareness for Multimodal Reasoning---An Interactive Agricultural Demonstration}
\titlerunning{MIRAGE-CARMA: An Interactive Agricultural Demonstration}

% TODO(author): Confirm author order, affiliations, and corresponding-author contact.
\author{Shashank Singh \and Tushar Gupta \and John F. Reid}
\authorrunning{S. Singh et al.}
\institute{University of Illinois Urbana--Champaign, Urbana, IL, USA}

\maketitle

\begin{abstract}
Agricultural questions grounded in images often cannot be resolved from visual
content alone: useful advice may depend on crop identity, geographic location,
season, and region-specific knowledge. We present MIRAGE-CARMA, an operational
prototype for context-adaptive evidence acquisition and multimodal response
generation. Given an image, question, and location, an evidence controller can
enrich the textual query, evaluate several applicable metadata-constrained
retrieval strategies together with unrestricted semantic retrieval, and select
the strategy providing the strongest relevant evidence. A confidence-routing
heuristic can escalate weak retrieval to domain-constrained You.com search;
acquired pages are ingested into a mutable, run-scoped Qdrant collection and
retrieved again without modifying the immutable curated collection. The
retrieved text is then supplied with the original image and question to a
vision--language model. The demonstration makes this process inspectable: users
can observe the effective query, strategy diagnostics, evidence provenance and
similarities, confidence, web-search decisions and URLs, runtime-memory writes,
and the final multimodal response. It is intended for IR researchers and
practitioners studying metadata-aware retrieval, adaptive RAG, and auditable
multimodal systems rather than as a replacement for professional agricultural
diagnosis.
\keywords{retrieval-augmented generation \and multimodal retrieval \and
metadata filtering \and adaptive retrieval \and agriculture}
\end{abstract}

\noindent\textbf{Demo availability:}
Live system: \url{DEMO_URL}\quad
Video: \url{VIDEO_URL}\quad
Source code: \url{CODE_URL}
% TODO(artifact): Replace at least DEMO_URL or VIDEO_URL before submission and
% replace CODE_URL if a public repository/release is made available.

\section{Introduction}

The agricultural MIRAGE benchmark frames assistance as expert-guided,
multimodal information seeking: a question can be accompanied by one or more
images and geographic context, while the desired response may require
identification, explanation, or management guidance~\cite{dongre2025mirage}.
An image can show visible symptoms, yet whether a recommendation is appropriate
may also depend on the crop, location, and season. This makes the task a useful
case for retrieval-augmented generation (RAG), which supplements model
parameters with inspectable and updateable external evidence~\cite{lewis2020rag}.
However, a fixed semantic search neither guarantees contextual fit nor provides
a response when the curated collection is insufficient.

MIRAGE-CARMA (Context-Adaptive Retrieval with Metadata Awareness for
Multimodal Reasoning) demonstrates an alternative: retrieval is treated as an
observable, adaptive process preceding multimodal answer generation. Its
\emph{progressive metadata filtering} evaluates all applicable strategies
formed from USDA plant hardiness zone, month--year, and title constraints, plus
unrestricted semantic retrieval, and selects the strongest strategy. This is
not a sequential fallback: every applicable candidate is evaluated under the
same selection rule. When
the selected evidence has low confidence, an evidence controller can search
external sources, add acquired passages to isolated runtime memory, retrieve
again, and reevaluate confidence. This corrective pattern is related to
adaptive and corrective RAG~\cite{asai2024selfrag,yan2024crag}, while the use of
structured filters connects it to metadata-aware RAG~\cite{poliakov2025multimetarag}.

The demonstration has three purposes: to show how agricultural context changes
retrieval, to expose why the system escalates beyond its curated knowledge, and
to make the provenance of the final VLM context inspectable. Its scope is
textual evidence acquisition for image-conditioned agricultural questions; the
evidence controller does not inspect the image or retrieve images. The intended
audience comprises IR researchers interested in metadata and adaptive
retrieval, multimodal-RAG developers, and agricultural-technology practitioners
who need to audit where supporting context came from.

\section{MIRAGE-CARMA System}

\paragraph{Inputs and query construction.}
The user supplies a question, one or more images, and a state or state--county
location. Location is retained in the effective query and, when resolvable, is
mapped to a USDA plant hardiness zone---a climate proxy based on average annual
extreme minimum temperature, not a complete model of local agricultural
conditions~\cite{usda2023phzm}. An optional crop-enrichment step may insert an
implicit crop name from a state-specific allowlist. The insertion is accepted
only if the original question remains a character-level subsequence; otherwise
the original text is used.

\paragraph{Progressive metadata filtering.}
Questions and passages are embedded with
\texttt{BAAI/bge-base-en-v1.5}. For each collection, the retriever constructs
every strategy supported by the available zone, \texttt{month\_year}, and
\texttt{title} values, and always includes semantic-only search. Each candidate
performs cosine search in Qdrant. Strategies with sufficient semantically
relevant hits are compared using the mean transformed score
$\frac{1}{n}\sum_i(2-s_i)^{-1}$, where $s_i$ is the raw Qdrant cosine score;
the maximum-scoring strategy is selected. Thus metadata constrains the search
space when it helps, while semantic-only retrieval remains a competing
strategy when metadata is missing or unhelpful.

Retrieval runs independently over two stores. \texttt{mirage\_base\_build} is
the immutable curated Qdrant collection. A second, mutable Qdrant collection is
created for the run. The winning result lists from the two collections are
merged, deduplicated by content hash (with base evidence preferred for a
duplicate), globally ordered by raw cosine similarity, and truncated to the
evidence budget. Each passage retains whether it came from the curated or
runtime collection.

\paragraph{Confidence-aware acquisition and generation.}
A fixed routing heuristic combines the relevant passages' mean similarity,
count, score consistency, and the scope of the selected retrieval strategy. It
is an operational routing score, not a calibrated probability of correctness.
For low-confidence evidence, the controller is instructed to extract keywords,
query You.com, optionally constrain the query to location-linked academic
domains, ingest successfully acquired pages into the runtime collection, and
repeat retrieval and confidence evaluation. The policy is prompt-directed by a
tool-using model rather than encoded as a deterministic state machine; the demo
therefore displays the executed calls as well as the intended route.

The resulting textual passages are concatenated with the effective question
and sent, together with the user images, to the answer VLM. Runtime evidence can
be reused by later questions in the same run, but it never mutates the curated
collection. Figure~\ref{fig:workflow} summarizes the boundary between the two
stores and the two model roles.

\begin{figure}[t]
  \centering
  \fbox{%
    \begin{minipage}[c][3.65cm][c]{0.94\linewidth}
      \centering\small
      \textbf{Question + image(s) + location}
      $\rightarrow$ \textbf{query enrichment}
      $\rightarrow$ \textbf{progressive metadata filtering}\\[2mm]
      \begin{tabular}{c@{\hspace{8mm}}c}
      \fbox{\parbox{0.34\linewidth}{\centering Immutable curated Qdrant\\
      \texttt{mirage\_base\_build}}} &
      \fbox{\parbox{0.34\linewidth}{\centering Mutable run-scoped Qdrant\\
      retrieved and updated at runtime}}
      \end{tabular}\\[2mm]
      $\downarrow$ merge, deduplicate, rank $\downarrow$\\[-1mm]
      \textbf{confidence evaluation}
      $\rightarrow$ \emph{if low:} domain-constrained web search
      $\rightarrow$ runtime ingestion $\rightarrow$ retrieve again\\[1mm]
      \textbf{retrieval context + original image(s)}
      $\rightarrow$ \textbf{multimodal answer generation}
    \end{minipage}}
  \caption{MIRAGE-CARMA architecture and demonstration workflow. The figure
  placeholder reserves space for the final graphic. Web-acquired evidence is
  written only to the mutable run-scoped collection; the curated collection
  remains immutable.}
  \label{fig:workflow}
\end{figure}

\section{Demonstration Experience}

An attendee begins from a prepared agricultural case or supplies a question,
image, and location. The workspace keeps the input and answer visible while an
execution trace expands between them. This layout lets a user inspect not only
the returned response but also how context changed the evidence supplied to the
VLM. Table~\ref{tab:trace} lists the principal views.

\begin{table}[t]
\caption{Information exposed during one demonstration run.}
\label{tab:trace}
\centering
\small
\begin{tabular}{p{0.24\linewidth}p{0.69\linewidth}}
\hline
\textbf{View} & \textbf{What the attendee can inspect} \\
\hline
Input & Original image(s), question, and location beside the effective or
crop-enriched query. \\
Retrieval & Applicable metadata strategies and semantic-only retrieval,
selected strategy, filters, raw similarities, passage text, provenance, and
curated/runtime source. \\
Routing & Confidence score and category, its diagnostics, and whether low
confidence triggered external acquisition. \\
Web and memory & Search keywords, You.com result URLs, fetch/ingestion status,
new runtime passages, and the second retrieval after ingestion. \\
Generation & The exact retrieved context admitted to generation and the final
image-conditioned answer. \\
\hline
\end{tabular}
\end{table}

The central interaction is a controlled rerun. A visitor can modify the
location or question while retaining the image and compare the eligible
filters, selected strategy, and evidence. They can also inspect an
augmentation case before and after ingestion: the trace first shows weak or
insufficient curated evidence, then the external URLs and chunks added to
runtime memory, and finally the revised retrieval state. A subsequent related
query demonstrates reuse by identifying passages whose source is the runtime
collection. The interface distinguishes a configured capability from observed
behavior: for example, domain filtering being enabled is shown separately from
the search query and URLs that were actually produced.

The displayed trace is derived from structured retrieval state and live tool
events rather than from the evidence controller's free-text terminal message.
This preserves selected-strategy diagnostics and passage-level cosine scores,
and makes deviations from the prompt-directed acquisition policy visible. The
answer is presented as evidence-supported model output, not as a field
diagnosis or treatment prescription.
% TODO(demo-ui): Confirm that the deployed trace adapter persists and displays
% all live search URLs, individual ingestion outcomes, and the admitted final
% context; the batch JSONL summary alone does not retain every tool event.

\section{Implementation and Demonstration Scenarios}

The prototype uses Qdrant server mode for concurrent retrieval and ingestion.
The live curated and runtime collections use cosine distance; the Qdrant server
version verified during the research runs is 1.18.0. Payload indexes support
hardiness zone, month--year, title, and content hash. Web pages are searched
through You.com, extracted and chunked, deduplicated against both enabled
collections, embedded, and written only to runtime memory. The research batch
deployment uses four NVIDIA A40 GPUs at a time, separating evidence-controller
workers from final response generation. The demonstration retains the same
service boundaries, although its final public deployment and VLM checkpoint
must be recorded with the artifact release.
% TODO(deployment): Record the final public-demo VLM checkpoint, revision,
% serving configuration, and hardware; exact Qwen-3 assignments are not yet
% verified and are intentionally absent from this draft.

We will present three complementary scenarios. \emph{Contextual retrieval}
holds an image/question pair fixed while changing location, revealing how a
derived hardiness zone alters eligible strategies without making metadata
mandatory. \emph{Confidence-aware escalation} starts with weak internal
evidence, displays keyword extraction and domain-constrained search, and shows
the post-ingestion retrieval. \emph{Runtime reuse} issues a related follow-up
within the same run and identifies evidence retrieved from the mutable store.
Together, the scenarios cover the fast path, the adaptive path, and the
stateful-memory behavior most relevant to an IR audience.

The repository also supports one direct-generation baseline and eight
configured RAG variants (nine experimental conditions in total). These
configurations are useful for development and controlled study, but the demo is
organized around the end-to-end adaptive system rather than an ablation table.
The final preload corpus and crop-dictionary artifacts are still being
finalized, so this draft deliberately gives no corpus-size or crop-coverage
statistics. Likewise, it reports no effectiveness results or exact model
assignments that have not yet been finalized.

\section{Conclusion}

MIRAGE-CARMA turns adaptive evidence acquisition into the object of the
demonstration. Attendees can follow contextual inputs through progressive
metadata filtering, confidence-aware web augmentation, isolated runtime
memory, and multimodal response generation. The prototype exposes both the
evidence that supports an answer and the routing decisions that obtained it,
while keeping live acquisition separate from the immutable curated collection.
Its present scope is U.S.-oriented agricultural question answering: hardiness
zones are coarse geographic proxies, live search is time-dependent, and the
confidence score is heuristic. These constraints also motivate the demo's
emphasis on inspectability rather than unsupported claims of diagnostic
authority or finalized empirical gains.

\bibliographystyle{splncs04}
\bibliography{references}

\end{document}
